% chapters/11/latihan.tex
\bagian{Latihan dan Refleksi}

\begin{enumerate}
  \item Jelaskan pekerjaan yang dihasilkan lapisan keputusan dan lapisan
        navigasi. Mengapa keduanya perlu dipisahkan?

  \item Sebuah pencarian jalur memeriksa 150 petak. Apabila tiga NPC
        menghitung ulang jalurnya setiap bingkai pada 60 bingkai per detik,
        hitunglah banyaknya petak yang diperiksa setiap detik.

  \item Jelaskan dua sebab pencarian jalur perlu diulang, dan apa yang terjadi
        apabila salah satunya dilupakan.

  \item Jelaskan mengapa pengikut jalur memakai perulangan \code{while},
        bukan \code{if}. Sebutkan gejala yang muncul apabila keduanya tertukar.

  \item Jelaskan mengapa toleransi jarak pada pengikut jalur diperlukan.
        Apa yang terjadi apabila nilainya nol?

  \item Sebutkan tiga segi penilaian mutu perilaku NPC beserta ukuran
        masing-masing, lalu jelaskan mengapa ketiganya saling bertentangan.

  \item Pada Tabel~\ref{tab:evaluasi-selang}, hitunglah rasio beban dan selisih
        mutu antara selang 0,15 detik dan 0,60 detik. Berikan kesimpulan Anda
        berdasarkan kedua angka itu.

  \item Sebutkan empat syarat skenario uji yang berguna, lalu periksa apakah
        jalur pemain yang dipakai program bab ini memenuhi keempatnya.
\end{enumerate}

\begin{refleksibox}
\begin{enumerate}
  \item Pada langkah keenam praktikum, perbaikan mana yang Anda pilih, dan
        berapa besar perbaikannya menurut pengukuran? Apakah hasilnya sesuai
        dugaan kelompok Anda?
  \item Separuh nilai mini proyek berkaitan dengan pengujian
        dan penggabungan. Menurut pengalaman kelompok Anda, bagian mana yang
        paling banyak memakan waktu: menulis kode, menggabungkan, atau menguji?
  \item Blok Python telah selesai. Sebutkan satu hal yang pada awalnya terasa
        sulit, tetapi sekarang terasa jelas, dan jelaskan apa yang membuatnya
        menjadi jelas.
\end{enumerate}
\end{refleksibox}
